\section{Fault Tolerance, Recovery, and Long-Term Deployment}
\label{sec:reliability}

Fault detection, isolation, accommodation, and restoration need not occur together, since a controller may tolerate bias of unknown source (Table~\ref{S-tab:reliability} in the supplementary material).

\subsection{Fault Estimation and Tolerant Control}

Analytic redundancy provides diagnosis without duplicating every sensor. Kinematics-based detection compares position, angular velocity, and acceleration through transport relations, with faults injected into recorded rover data~\cite{P_E51}. A linear parameter-varying (LPV) vehicle diagnoser combines zonotopic set-membership estimation with invariant-set bounds, in which separation from the healthy set establishes detectability and overlap between fault classes prevents unique isolation~\cite{P_N16}. Tolerant controllers instead accommodate faults in specified channels, each relying on a different remaining capability. Mixed-gain prescribed-performance control tolerates partial wheel-effectiveness loss~\cite{P_M15}, whereas virtual-error control tolerates heading and wheel-speed faults but assumes exact position~\cite{P_M01}. Adaptive sliding-mode control handles injected actuator faults and saturation on industrial hardware as long as some actuation remains~\cite{P_E52}, while fixed-time adaptive tracking accommodates dead zones and modeled actuator faults in MATLAB studies~\cite{P_S15}. An observer-based event-triggered controller reconstructs modeled faults before issuing updates in hardware-in-the-loop tests of a low-speed, no-slip vehicle~\cite{P_F02}.

\subsection{Integrity and Failure Monitoring}

Monitors differ in the failure event they characterize. Reliable Monte Carlo localization (MCL) estimates a latent localization-success state alongside pose and triggers free-space global relocalization~\cite{P_M35}. Integrity analysis for extended Kalman filter (EKF) SLAM accumulates landmark fault risk into the integrity risk \eqref{eq:integrity}~\cite{P_X07}. In its recorded example, this risk crosses the threshold only after localization fails. A prior-line-map protection level predicts estimation error, with operational alerting left for future work~\cite{P_F01}. Monitors of learned policies use statistical alarm definitions, as conformal anomaly detection controls the nominal false-alarm rate under exchangeability~\cite{P_E54}, while Sentinel pairs fast action-consistency monitoring with slower vision-language progress assessment~\cite{P_F07}. PAAD predicts anomalies along a planned path on field recordings in which people reposition the robot after failures~\cite{P_N18}. RADIUM moves failure prediction before deployment by sampling environment parameters from a distribution conditioned on failure and repairs the policy on the discovered cases, reporting up to five times higher robustness for a repaired vision-in-the-loop policy on a 1/10-scale car~\cite{P_R36}.

\subsection{From Diagnosis to an Executed Response}

The response must match the location of the fault in the stack. In the estimator, fault-tolerant prior-map collaboration excludes faulty observations, inflates faulty-landmark covariance, and specifies reinitialization for odometry faults without evaluating an executed reset~\cite{P_F03}. At the planning level, a GP predictor of safe-corridor optimizer failure triggers a stop--reposition--resume sequence~\cite{P_E59}. With the same planner in simulation, recovery yields eight successes in ten runs against three in ten in a dense world with 1\,m doorways, whereas both variants succeed in all ten runs in a sparse world with 2\,m doorways. However, the predictor posterior concerns planner failure rather than collision. At the level of the controller and supervisor, a fault-aware CBF controller optimizes over candidate sensor or actuator fault sets in numerical examples~\cite{P_E55}, whereas FORTRESS caches slow semantic reasoning for runtime reach-avoid fallback on a quadrotor~\cite{P_E60}.

Implemented recovery policies differ in what they restore, from localization to task feasibility. A nonlinear MPC (NMPC) visual navigation system stops after tracking loss and attempts relocalization~\cite{P_N10}, whereas behavior-tree supervision on CENTAURO chooses among contact repositioning, reversing, replanning, and human assistance~\cite{P_N17}. Task-level policies change the mission instead, as resilient temporal-logic planning relaxes task windows or removes tasks when local correction fails~\cite{P_N15} and predictive swarm fault tolerance withdraws an affected robot to maintenance before complete failure~\cite{P_N13}. Local recovery acts on the trajectory or the observations, through hybrid-A* reinitialization of an infeasible local trajectory in RTEB~\cite{P_S16}, return toward familiar visual states in SCALE~\cite{P_S38}, and repulsive corrections from estimated depth to a visual policy in CARE~\cite{P_S47}. However, neither familiarity nor estimated depth supplies a clearance certificate. Help Me Through learns viewpoint actions from about 50\,h of human demonstrations, or 20\,h followed by online fine-tuning, to keep monocular SLAM tracking through texture-poor transitions, with a physical test that runs the policy on live images while a person teleoperates path planning~\cite{P_R33}. When a robot cannot restore a capability, the team can reassign its task, as in energy-aware allocation under staged capability loss on Robotarium~\cite{P_S17} and in R-FAC, whose log-sum-exp value factorization keeps search policies effective when members fail without notification, assuming perfect target detection~\cite{P_R35}.

\subsection{Resources and Sustained Operation}

Resource choices determine whether a response remains available. CODEI selects bodies, sensors, algorithms, and computing units from catalogs by coupling simulated planner queries to characterized perception errors~\cite{P_M11}. Joint runtime control of motor speed and processor voltage and frequency on a wheeled rover reduces total energy by up to 36.34\% relative to prior methods, with measurements showing that computation can dominate the energy of slow motion~\cite{P_R32}. Fleet availability adds charging constraints, as in cyclic scheduling of operating and charging periods under a station budget~\cite{P_E58} and in persistent-monitoring routes that minimize team size under revisit-latency and recharging limits with deterministic discharge~\cite{P_R34}. A ten-home assistive deployment combines docking rules, cloud orchestration, and remote supervision, with internet problems and localization failures among the causes of interruption~\cite{P_M37}. A cloud--edge system offloads mapping and trajectory generation over 5G while keeping MPC tracking and collision avoidance on Jetson-equipped Ackermann robots, falling back to a leader-based local network~\cite{P_R31}.

Maintenance and learning form a second operational loop. Predictive data acquisition selects new training observations for teach-repeat-learn navigation from visual matching ambiguity~\cite{P_N06}. SELFI stabilizes deployment-time reinforcement learning with a model-based objective and counts interventions during learning~\cite{P_S14}, while LiReN combines foundation-model adaptation with pseudo-resets, keepout rules, charging, and help escalation~\cite{P_S41}. Counterfactual replay selects among adapting policies at the cost of between-trial computation~\cite{P_S18}.

\subsection{Synthesis}

Recovery quality in long-term service concerns both restored operation and the interruptions needed to achieve it, which depend on charging, energy, network access, offboard computation, and human assistance. The studies reviewed here contain \RelIntegrityCount{} of the \IntegrityCount{} integrity or failure-risk assessments in the sample (Table~\ref{tab:evaluation}), mostly with injected or simulated faults.
